% Corrected French for printed p.139.
% Corrections: nue -> nul; restore the missing middle term B in the
% second exact sequence. A provisional A' -> A'' correction was reversed
% after high-detail inspection showed that the authority already reads A''.
\SourcePage{139}

\LayerTarget{tohoku:def:quotient-composition:p2}
\noindent sous-truc $B''$ de $B$ dans le quotient $B'$, $Q$ est aussi
isomorphe à la coimage de ce morphisme, et est donc à la fois un sous-truc de
$B'$ et un truc-quotient de $B''$. En diminuant au besoin le sous-truc $A'$ de
$A$ et le quotient $C'$ de $C$, on peut supposer que $u$ et $v$ proviennent
respectivement de morphismes (notés encore par la même lettre)
$A'\longrightarrow Q$ et $Q\longrightarrow C'$. On peut donc prendre le
composé $vu\in\mathrm{Hom}_C(A',C')$, et on vérifie que l'élément de
$\mathrm{Hom}_{C/C'}(A,C)$ qu'il définit ne dépend que de $u'$ et $v'$, on le
notera $v'u'$. Il n'y a aucune difficulté à prouver que la loi de composition
ainsi définie est bilinéaire, associative, et que la classe dans
$\mathrm{Hom}_{C/C'}(A,A)$ de l'application identique $I_A$ est une unité
universelle, donc que $C/C'$ est une catégorie additive ; et enfin que c'est
même une catégorie abélienne. Nous ne ferons pas ici cette vérification
extrêmement fastidieuse.

\LayerTarget{tohoku:prop:quotient-category-properties}
Ainsi \textit{$C/C'$ apparaît comme une catégorie abélienne, de plus le
foncteur identique $F$ de $C$ dans $C/C'$ est exact} (et permute en
particulier aux noyaux, conoyaux, images et coimages), \textit{$F(A)$ est nul
si et seulement si $A\in C'$, enfin tout objet de $C/C'$ est de la forme
$F(A)$, $A\in C$.} Ce sont ces faits (qui essentiellement caractérisent la
catégorie quotient) qui permettent d'appliquer le langage
\textquotedblleft mod. $C'$\textquotedblright\ en toute sécurité, ce langage
signifiant simplement qu'on se place dans la catégorie abélienne quotient.

\LayerTarget{tohoku:remark:quotient-spectral-sequences}
Il est particulièrement commode pour exploiter, quand on a une suite spectrale
(cf. 2.4) dans $C$, le fait que certains termes de cette suite spectrale sont
dans $C'$ : réduisant mod. $C'$ (i. e. appliquant le foncteur $F$) on trouve
une suite spectrale dans $C/C'$ où les termes correspondants sont nuls, d'où
souvent des suites exactes mod. $C'$ à l'aide des critères usuels d'obtention
de suites exactes à partir de suites spectrales dont certains termes sont
nuls.

\LayerTarget{tohoku:chapter:2}
\bigskip
\begin{center}
{\large\bfseries Chapitre II Algèbre Homologique dans les catégories
abéliennes.}
\end{center}

\LayerTarget{tohoku:section:2.1}
\noindent\textbf{2.1. $\partial$-foncteurs et
$\partial^*$-foncteurs.}

\LayerTarget{tohoku:def:delta-functor:p1}
Soient $C$ une catégorie abélienne, $C'$ une catégorie additive, $a$ et $b$
deux entiers (pouvant être égaux à $+\infty$ ou $-\infty$) tels que
$a+1<b$. Un \textit{$\partial$-foncteur covariant de $C$ dans $C'$, à degrés
$a<i<b$}, est un système $T=(T^i)$ de foncteurs covariants additifs de $C$
dans $C'$, $(a<i<b)$, plus la donnée, pour tout $i$ tel que $a<i<b-1$, et
toute suite exacte
$0\longrightarrow A'\longrightarrow A\longrightarrow A''\longrightarrow 0$
dans $C$, d'un morphisme

\[
 \partial:T^i(A'')\longrightarrow T^{i+1}(A')
\]

\noindent (\textquotedblleft l'homomorphisme bord\textquotedblright\ ou
\textquotedblleft connecting homomorphism\textquotedblright), les axiomes
suivants étant supposés satisfaits :

\LayerTarget{tohoku:axiom:delta-functor:1}
\noindent (i) Si on a une deuxième suite exacte
$0\longrightarrow B'\longrightarrow B\longrightarrow B''\longrightarrow 0$ et
un homomorphisme de la première suite exacte dans la seconde, alors les
diagrammes correspondants

\LayerTarget{tohoku:diagram:delta-naturality}
\[
\begin{array}{ccc}
T^i(A'') & \mathop{\longrightarrow}\limits^{\partial} & T^{i+1}(A')\\
\big\downarrow && \big\downarrow\\
T^i(B'') & \mathop{\longrightarrow}\limits^{\partial} & T^{i+1}(B')
\end{array}
\]

\noindent sont commutatifs.

\LayerTarget{tohoku:axiom:delta-functor:2:p1}
\noindent (ii) Pour toute suite exacte
$0\longrightarrow A'\longrightarrow A\longrightarrow A''\longrightarrow 0$,
la suite de morphismes associée

\TohokuContinuation{140}
